% Options for packages loaded elsewhere
% Options for packages loaded elsewhere
\PassOptionsToPackage{unicode}{hyperref}
\PassOptionsToPackage{hyphens}{url}
\PassOptionsToPackage{dvipsnames,svgnames,x11names}{xcolor}
%
\documentclass[
  letterpaper,
  DIV=11,
  numbers=noendperiod]{scrreprt}
\usepackage{xcolor}
\usepackage[heightrounded]{geometry}
\usepackage{amsmath,amssymb}
\setcounter{secnumdepth}{-\maxdimen} % remove section numbering
\usepackage{iftex}
\ifPDFTeX
  \usepackage[T1]{fontenc}
  \usepackage[utf8]{inputenc}
  \usepackage{textcomp} % provide euro and other symbols
\else % if luatex or xetex
  \usepackage{unicode-math} % this also loads fontspec
  \defaultfontfeatures{Scale=MatchLowercase}
  \defaultfontfeatures[\rmfamily]{Ligatures=TeX,Scale=1}
\fi
\usepackage{lmodern}
\ifPDFTeX\else
  % xetex/luatex font selection
\fi
% Use upquote if available, for straight quotes in verbatim environments
\IfFileExists{upquote.sty}{\usepackage{upquote}}{}
\IfFileExists{microtype.sty}{% use microtype if available
  \usepackage[]{microtype}
  \UseMicrotypeSet[protrusion]{basicmath} % disable protrusion for tt fonts
}{}
\makeatletter
\@ifundefined{KOMAClassName}{% if non-KOMA class
  \IfFileExists{parskip.sty}{%
    \usepackage{parskip}
  }{% else
    \setlength{\parindent}{0pt}
    \setlength{\parskip}{6pt plus 2pt minus 1pt}}
}{% if KOMA class
  \KOMAoptions{parskip=half}}
\makeatother
% Make \paragraph and \subparagraph free-standing
\makeatletter
\ifx\paragraph\undefined\else
  \let\oldparagraph\paragraph
  \renewcommand{\paragraph}{
    \@ifstar
      \xxxParagraphStar
      \xxxParagraphNoStar
  }
  \newcommand{\xxxParagraphStar}[1]{\oldparagraph*{#1}\mbox{}}
  \newcommand{\xxxParagraphNoStar}[1]{\oldparagraph{#1}\mbox{}}
\fi
\ifx\subparagraph\undefined\else
  \let\oldsubparagraph\subparagraph
  \renewcommand{\subparagraph}{
    \@ifstar
      \xxxSubParagraphStar
      \xxxSubParagraphNoStar
  }
  \newcommand{\xxxSubParagraphStar}[1]{\oldsubparagraph*{#1}\mbox{}}
  \newcommand{\xxxSubParagraphNoStar}[1]{\oldsubparagraph{#1}\mbox{}}
\fi
\makeatother

\usepackage{color}
\usepackage{fancyvrb}
\newcommand{\VerbBar}{|}
\newcommand{\VERB}{\Verb[commandchars=\\\{\}]}
\DefineVerbatimEnvironment{Highlighting}{Verbatim}{commandchars=\\\{\}}
% Add ',fontsize=\small' for more characters per line
\newenvironment{Shaded}{}{}
\newcommand{\AlertTok}[1]{\textcolor[rgb]{1.00,0.33,0.33}{\textbf{#1}}}
\newcommand{\AnnotationTok}[1]{\textcolor[rgb]{0.42,0.45,0.49}{#1}}
\newcommand{\AttributeTok}[1]{\textcolor[rgb]{0.84,0.23,0.29}{#1}}
\newcommand{\BaseNTok}[1]{\textcolor[rgb]{0.00,0.36,0.77}{#1}}
\newcommand{\BuiltInTok}[1]{\textcolor[rgb]{0.84,0.23,0.29}{#1}}
\newcommand{\CharTok}[1]{\textcolor[rgb]{0.01,0.18,0.38}{#1}}
\newcommand{\CommentTok}[1]{\textcolor[rgb]{0.42,0.45,0.49}{#1}}
\newcommand{\CommentVarTok}[1]{\textcolor[rgb]{0.42,0.45,0.49}{#1}}
\newcommand{\ConstantTok}[1]{\textcolor[rgb]{0.00,0.36,0.77}{#1}}
\newcommand{\ControlFlowTok}[1]{\textcolor[rgb]{0.84,0.23,0.29}{#1}}
\newcommand{\DataTypeTok}[1]{\textcolor[rgb]{0.84,0.23,0.29}{#1}}
\newcommand{\DecValTok}[1]{\textcolor[rgb]{0.00,0.36,0.77}{#1}}
\newcommand{\DocumentationTok}[1]{\textcolor[rgb]{0.42,0.45,0.49}{#1}}
\newcommand{\ErrorTok}[1]{\textcolor[rgb]{1.00,0.33,0.33}{\underline{#1}}}
\newcommand{\ExtensionTok}[1]{\textcolor[rgb]{0.84,0.23,0.29}{\textbf{#1}}}
\newcommand{\FloatTok}[1]{\textcolor[rgb]{0.00,0.36,0.77}{#1}}
\newcommand{\FunctionTok}[1]{\textcolor[rgb]{0.44,0.26,0.76}{#1}}
\newcommand{\ImportTok}[1]{\textcolor[rgb]{0.01,0.18,0.38}{#1}}
\newcommand{\InformationTok}[1]{\textcolor[rgb]{0.42,0.45,0.49}{#1}}
\newcommand{\KeywordTok}[1]{\textcolor[rgb]{0.84,0.23,0.29}{#1}}
\newcommand{\NormalTok}[1]{\textcolor[rgb]{0.14,0.16,0.18}{#1}}
\newcommand{\OperatorTok}[1]{\textcolor[rgb]{0.14,0.16,0.18}{#1}}
\newcommand{\OtherTok}[1]{\textcolor[rgb]{0.44,0.26,0.76}{#1}}
\newcommand{\PreprocessorTok}[1]{\textcolor[rgb]{0.84,0.23,0.29}{#1}}
\newcommand{\RegionMarkerTok}[1]{\textcolor[rgb]{0.42,0.45,0.49}{#1}}
\newcommand{\SpecialCharTok}[1]{\textcolor[rgb]{0.00,0.36,0.77}{#1}}
\newcommand{\SpecialStringTok}[1]{\textcolor[rgb]{0.01,0.18,0.38}{#1}}
\newcommand{\StringTok}[1]{\textcolor[rgb]{0.01,0.18,0.38}{#1}}
\newcommand{\VariableTok}[1]{\textcolor[rgb]{0.89,0.38,0.04}{#1}}
\newcommand{\VerbatimStringTok}[1]{\textcolor[rgb]{0.01,0.18,0.38}{#1}}
\newcommand{\WarningTok}[1]{\textcolor[rgb]{1.00,0.33,0.33}{#1}}

\usepackage{longtable,booktabs,array}
\usepackage{calc} % for calculating minipage widths
% Correct order of tables after \paragraph or \subparagraph
\usepackage{etoolbox}
\makeatletter
\patchcmd\longtable{\par}{\if@noskipsec\mbox{}\fi\par}{}{}
\makeatother
% Allow footnotes in longtable head/foot
\IfFileExists{footnotehyper.sty}{\usepackage{footnotehyper}}{\usepackage{footnote}}
\makesavenoteenv{longtable}
\usepackage{graphicx}
\makeatletter
\newsavebox\pandoc@box
\newcommand*\pandocbounded[1]{% scales image to fit in text height/width
  \sbox\pandoc@box{#1}%
  \Gscale@div\@tempa{\textheight}{\dimexpr\ht\pandoc@box+\dp\pandoc@box\relax}%
  \Gscale@div\@tempb{\linewidth}{\wd\pandoc@box}%
  \ifdim\@tempb\p@<\@tempa\p@\let\@tempa\@tempb\fi% select the smaller of both
  \ifdim\@tempa\p@<\p@\scalebox{\@tempa}{\usebox\pandoc@box}%
  \else\usebox{\pandoc@box}%
  \fi%
}
% Set default figure placement to htbp
\def\fps@figure{htbp}
\makeatother





\setlength{\emergencystretch}{3em} % prevent overfull lines

\providecommand{\tightlist}{%
  \setlength{\itemsep}{0pt}\setlength{\parskip}{0pt}}



 


\usepackage{fvextra}
\DefineVerbatimEnvironment{Highlighting}{Verbatim}{breaklines,commandchars=\\\{\}}
\KOMAoption{captions}{tableheading}
\makeatletter
\@ifpackageloaded{tcolorbox}{}{\usepackage[skins,breakable]{tcolorbox}}
\@ifpackageloaded{fontawesome5}{}{\usepackage{fontawesome5}}
\definecolor{quarto-callout-color}{HTML}{909090}
\definecolor{quarto-callout-note-color}{HTML}{0758E5}
\definecolor{quarto-callout-important-color}{HTML}{CC1914}
\definecolor{quarto-callout-warning-color}{HTML}{EB9113}
\definecolor{quarto-callout-tip-color}{HTML}{00A047}
\definecolor{quarto-callout-caution-color}{HTML}{FC5300}
\definecolor{quarto-callout-color-frame}{HTML}{acacac}
\definecolor{quarto-callout-note-color-frame}{HTML}{4582ec}
\definecolor{quarto-callout-important-color-frame}{HTML}{d9534f}
\definecolor{quarto-callout-warning-color-frame}{HTML}{f0ad4e}
\definecolor{quarto-callout-tip-color-frame}{HTML}{02b875}
\definecolor{quarto-callout-caution-color-frame}{HTML}{fd7e14}
\makeatother
\makeatletter
\@ifpackageloaded{caption}{}{\usepackage{caption}}
\AtBeginDocument{%
\ifdefined\contentsname
  \renewcommand*\contentsname{Table of contents}
\else
  \newcommand\contentsname{Table of contents}
\fi
\ifdefined\listfigurename
  \renewcommand*\listfigurename{List of Figures}
\else
  \newcommand\listfigurename{List of Figures}
\fi
\ifdefined\listtablename
  \renewcommand*\listtablename{List of Tables}
\else
  \newcommand\listtablename{List of Tables}
\fi
\ifdefined\figurename
  \renewcommand*\figurename{Figure}
\else
  \newcommand\figurename{Figure}
\fi
\ifdefined\tablename
  \renewcommand*\tablename{Table}
\else
  \newcommand\tablename{Table}
\fi
}
\@ifpackageloaded{float}{}{\usepackage{float}}
\floatstyle{ruled}
\@ifundefined{c@chapter}{\newfloat{codelisting}{h}{lop}}{\newfloat{codelisting}{h}{lop}[chapter]}
\floatname{codelisting}{Listing}
\newcommand*\listoflistings{\listof{codelisting}{List of Listings}}
\makeatother
\makeatletter
\makeatother
\makeatletter
\@ifpackageloaded{caption}{}{\usepackage{caption}}
\@ifpackageloaded{subcaption}{}{\usepackage{subcaption}}
\makeatother
\makeatletter
\@ifpackageloaded{tcolorbox}{}{\usepackage[skins,breakable]{tcolorbox}}
\makeatother
\makeatletter
\@ifundefined{shadecolor}{\definecolor{shadecolor}{HTML}{31BAE9}}{}
\makeatother
\makeatletter
\makeatother
\makeatletter
\ifdefined\Shaded\renewenvironment{Shaded}{\begin{tcolorbox}[boxrule=0pt, borderline west={3pt}{0pt}{shadecolor}, interior hidden, frame hidden, sharp corners, enhanced, breakable]}{\end{tcolorbox}}\fi
\makeatother
\usepackage{bookmark}
\IfFileExists{xurl.sty}{\usepackage{xurl}}{} % add URL line breaks if available
\urlstyle{same}
\hypersetup{
  colorlinks=true,
  linkcolor={blue},
  filecolor={Maroon},
  citecolor={Blue},
  urlcolor={Blue},
  pdfcreator={LaTeX via pandoc}}


\author{}
\date{}
\begin{document}

\RecustomVerbatimEnvironment{verbatim}{Verbatim}{
  showspaces = false,
  showtabs = false,
  breaksymbolleft={},
  breaklines
}

\renewcommand*\contentsname{Table of contents}
{
\hypersetup{linkcolor=}
\setcounter{tocdepth}{2}
\tableofcontents
}

\chapter{Using LLMs to understand and check
code}\label{using-llms-to-understand-and-check-code}

Tools like ChatGPT, Claude, Copilot or Gemini can write bash code in
seconds, so why would you still want to learn how to use bash?

For this, let's first understand what these tools are. Very briefly, all
these tools are built on a \textbf{large language model (LLM)}: a
program that has been trained on huge amounts of text, including a lot
of code, and that writes an answer by predicting which words are likely
to come next. On this page, I use ``LLM'' for all of these chat tools.

LLMs are useful, and some of you may already use them. This page is
about how to use them without losing control over your analysis. It is
written based on my own experience using LLMs, for people who work with
biological data and who have just started with the command line. You can
read it after the workshop, at your own pace.

The short version:

\begin{itemize}
\tightlist
\item
  \textbf{Learn the basics first}, with this tutorial, and don't let an
  LLM write code for you while you do. The basics give you the tools you
  need for all the other advice on this page
\item
  Both while you learn and later, use an LLM as a \textbf{tutor} that
  explains code, more than as a machine that writes code for you
\item
  Treat every piece of code an LLM writes as a \textbf{first draft}: run
  it on one sample and check the output before you use it for real
\item
  You need to be able to \textbf{explain} and \textbf{justify} every
  step of your analysis, whoever wrote the code
\item
  You need to \textbf{understand} what each option of a tool does to
  your data. For the code around it, you can \textbf{check} that it does
  what you want with the simple commands from day 1
\end{itemize}

\section{Why learn the command line if an LLM can write the
code?}\label{why-learn-the-command-line-if-an-llm-can-write-the-code}

This is a fair question. In my opinion, even if you use an LLM, someone
still has to:

\begin{itemize}
\tightlist
\item
  \textbf{Read the code}, to know what it will do to your files before
  you run it. This means understanding details like how exactly the data
  is summarized or how the options of a command, for example fastp,
  affect your output
\item
  \textbf{Run the code} in the right place: on the right computer, in
  the right folder, with the right paths, on the login node or as a job.
  The LLM does not know your folder structure, your file names or the
  HPC you work on
\item
  \textbf{Check the input and output}: what does your input data look
  like, did every sample get processed, how many reads are left, does
  that make sense for your data?
\item
  \textbf{Explain and justify} your choices: to your supervisor, to your
  co-authors, in the methods section of a paper. ``The LLM chose this
  cutoff'' is not really an acceptable answer if someone asks you about
  your choices
\end{itemize}

All four need the skills from this workshop. An LLM can help with each
of them, but it cannot take over the responsibility. That stays with
you.

\section{Start with the basics, without an
LLM}\label{start-with-the-basics-without-an-llm}

All the advice on this page depends on the basics from this tutorial. To
check code you need to be able to navigate the file system with the
terminal and know about basic commands like \texttt{ls}, \texttt{wc\ -l}
or \texttt{head}. To notice that a loop skipped a sample, you need to
have written a few loops yourself. And to see that an output looks
wrong, you need to know what a normal output looks like.

So while you learn the basics, with this tutorial and your first own
analyses, my advice is:

\begin{itemize}
\tightlist
\item
  \textbf{Don't use an LLM to write code.} Type the commands yourself,
  make mistakes and read the error messages. This is how you learn what
  normal looks like, and you need exactly that later, to spot when
  something is wrong. Struggle is a normal part of the learning
  experience and you don't want to outsource this to an LLM.
\item
  \textbf{Use an LLM only when you don't understand something and can't
  ask anyone}: not your supervisor, not a colleague, not the instructor.
  Even then, ask it to explain, not to give you the solution
\end{itemize}

Once you are comfortable with the basics, the rest of this page applies.

\section{Which tool can you use?}\label{which-tool-can-you-use}

\subsection{Rules at the UvA}\label{rules-at-the-uva}

Universities have their own rules for LLMs, mostly about privacy. At the
UvA (checked in October 2026, rules may change):

\begin{itemize}
\tightlist
\item
  The UvA offers its own tool,
  \href{https://www.uva.nl/en/about-the-uva/about-the-university/ai/ai-in-education/uva-ai-chat/uva-ai-chat.html}{UvA
  AI Chat}, to students and staff. It is free with your UvAnetID and
  works much like ChatGPT. According to UvA, what you type stays within
  the UvA environment and is not used to train the models
\item
  For \textbf{students}, UvA writes that commercial tools such as
  ChatGPT are
  \href{https://student.uva.nl/en/information/ai-tools-and-your-studies}{not
  allowed in the context of your studies}, because they don't meet the
  UvA privacy rules. Your study programme and your course decide what
  you may use an LLM for, and you are expected to be open about how you
  used it
\item
  For \textbf{everyone}, also in UvA AI Chat, the normal rules for
  privacy and research data still apply. The UvA's
  \href{https://www.uva.nl/en/about-the-uva/about-the-university/ai/ai-in-education/uva-ai-chat/uva-ai-chat-rai.html}{page
  on responsible use} says, for example, never to upload lists of names
  of students, patients or research participants, and to contact your
  faculty's data steward when in doubt
\end{itemize}

If you are not at the UvA, check the rules of your own institution.

\begin{tcolorbox}[enhanced jigsaw, colframe=quarto-callout-important-color-frame, opacitybacktitle=0.6, colbacktitle=quarto-callout-important-color!10!white, title=\textcolor{quarto-callout-important-color}{\faExclamation}\hspace{0.5em}{Important: Don't paste your data into an LLM}, rightrule=.15mm, coltitle=black, bottomtitle=1mm, opacityback=0, toptitle=1mm, titlerule=0mm, arc=.35mm, toprule=.15mm, bottomrule=.15mm, left=2mm, colback=white, leftrule=.75mm, breakable]

Don't share confidential information, unpublished data, sequences or
results with an LLM. You rarely need to. To ask about a problem, a few
lines are enough: the command you ran, the error message, and the first
lines of a file (\texttt{head\ -n\ 8} shows the first two reads of a
fastq file). Replace sample names or other details if they are
sensitive.

\end{tcolorbox}

\subsection{Different tools, same
advice}\label{different-tools-same-advice}

LLM tools differ a lot: free or paid, a chat window in your browser, or
a tool inside a code editor that can read your files and run commands
itself. Names, features and prices change every few months, so I don't
recommend a specific one here.

Everything on this page works with a free chat tool, such as UvA AI
Chat: you paste code, you ask a question, you read the answer. If you
use a tool that runs commands for you, the advice matters even more. You
see less of what happens, so checking the result becomes even more
important.

\section{How I suggest you use an
LLM}\label{how-i-suggest-you-use-an-llm}

These are the rules I give the students I supervise for how to interact
with an LLM. As I mentioned above, I advise students to not let an LLM
write their code and not outsource their learning. Your internships are
not only for you to produce results, they give you the time and space to
learn about the process of producing good and reliable results. However,
I am also aware that everyone has access to these tools, so these
recommendations also cover how I would use LLMs for coding:

\begin{itemize}
\tightlist
\item
  \textbf{Ask it to explain, not only to write.} When you learn a new
  concept, you can even tell it not to write code, for example: ``Give
  advice and best practices, but do NOT give me a code solution''
\item
  \textbf{Ask specific questions.} Give it the actual command, the
  actual error message and the biological context. The more specific the
  question, the more useful the answer
\item
  \textbf{Ask for the simplest version.} LLMs tend to write more code
  than you need: error handling for every possible case, a variable for
  every path. The code is not always wrong, but it is harder to
  understand what it does. Ask for the simplest version first, and add
  complexity only when you understand why you need it
\item
  \textbf{Treat generated code as a first draft.} Run it on one sample,
  look at the output, and check that it makes biological sense before
  you run it on all samples
\item
  \textbf{Don't let it rewrite code that fails.} This usually replaces
  code you don't understand with more code you don't understand. Ask it
  to explain the error or a specific line instead (more on this
  \hyperref[when-your-code-does-not-work]{below})
\item
  \textbf{Check suggested options against the help page of the version
  you use.} An LLM has learned from documentation of many versions of a
  tool, often older ones. Remember from day 2: Crunchomics has fastp
  1.0.1, and we installed 1.4.0. Options can change between versions.
  \texttt{fastp\ -\/-help} on your own installation is always more
  reliable than an LLM's suggestion
\item
  \textbf{Write down what each result means, in your own words.} Your
  notebook is not only a record of what you did. It is also a place to
  slow down and think. After each analysis step, write at least one
  sentence about what the result means for your data, for example:
  ``fastp removed 3\% of the reads'' and also think about whether that
  number is something that you expected. Make sure to write this
  yourself, not with an LLM
\end{itemize}

And the most important rule: \textbf{if you can't explain a step, don't
use it yet.} It is completely fine to ask for help with a step you don't
understand, from a person or from an LLM. Struggling with code is a
normal part of learning.

\begin{tcolorbox}[enhanced jigsaw, colframe=quarto-callout-color-frame, opacitybacktitle=0.6, colbacktitle=quarto-callout-color!10!white, title=\textcolor{quarto-callout-color}{\faInfo}\hspace{0.5em}{Exercise}, rightrule=.15mm, coltitle=black, bottomtitle=1mm, opacityback=0, toptitle=1mm, titlerule=0mm, arc=.35mm, toprule=.15mm, bottomrule=.15mm, left=2mm, colback=white, leftrule=.75mm, breakable]

On day 2, we ran fastp without any filtering options. Let's use an LLM
to find out whether it filtered all reads anyway:

\begin{enumerate}
\def\labelenumi{\arabic{enumi}.}
\tightlist
\item
  Ask in UvA AI Chat (or another chat tool that you are allowed to use):
  ``What does fastp do by default to filter reads? Which options control
  this, and what are their default values?'' Write down its answer
\item
  On Crunchomics, activate the \texttt{fastp\_1.4.0} environment and
  look at the help page. Find the options the LLM mentioned. Do the
  descriptions and default values match?
\item
  Deactivate the environment and look at the help page of the fastp
  1.0.1 of Crunchomics. Is there anything different?
\end{enumerate}

\begin{quote}
\textbf{Click me to see an answer}

\textbf{The help page of fastp 1.4.0 (question 2)}

\begin{Shaded}
\begin{Highlighting}[]
\CommentTok{\# Question 2: the help page of fastp 1.4.0}
\ExtensionTok{conda}\NormalTok{ activate fastp\_1.4.0}
\ExtensionTok{fastp} \AttributeTok{{-}{-}version}
\ExtensionTok{fastp} \AttributeTok{{-}{-}help} \DecValTok{2}\OperatorTok{\textgreater{}\&}\DecValTok{1} \KeywordTok{|} \FunctionTok{less}

\CommentTok{\# Question 3: the help page of the fastp of Crunchomics}
\ExtensionTok{conda}\NormalTok{ deactivate}
\ExtensionTok{fastp} \AttributeTok{{-}{-}version}
\ExtensionTok{fastp} \AttributeTok{{-}{-}help} \DecValTok{2}\OperatorTok{\textgreater{}\&}\DecValTok{1} \KeywordTok{|} \FunctionTok{less}
\end{Highlighting}
\end{Shaded}

\texttt{2\textgreater{}\&1} sends the error output to the same place as
the normal output. Some tools print their help page as error output, and
then \texttt{less} would not get it. Inside \texttt{less}, type
\texttt{/quality} and press Enter to jump to the word ``quality'', and
\texttt{n} to jump to the next match.

Look for the options about quality and length filtering. For fastp 1.4.0
on Crunchomics, these are the lines:

\begin{verbatim}
  -Q, --disable_quality_filtering      quality filtering is enabled by default. If this option is specified, quality filtering is disabled
  -q, --qualified_quality_phred        the quality value that a base is qualified. Default 15 means phred quality >=Q15 is qualified. (int [=15])
  -u, --unqualified_percent_limit      how many percents of bases are allowed to be unqualified (0~100). Default 40 means 40% (int [=40])
  -n, --n_base_limit                   if one read's number of N base is >n_base_limit, then this read/pair is discarded. Default is 5 (int [=5])
  -e, --average_qual                   if one read's average quality score <avg_qual, then this read/pair is discarded. Default 0 means no requirement (int [=0])
  -L, --disable_length_filtering       length filtering is enabled by default. If this option is specified, length filtering is disabled
  -l, --length_required                reads shorter than length_required will be discarded, default is 15. (int [=15])
\end{verbatim}

In plain words:

\begin{itemize}
\tightlist
\item
  \texttt{-q} and \texttt{-u}: a base with a quality score of 15 or more
  counts as good. A read is removed if more than 40\% of its bases are
  not good
\item
  \texttt{-n}: a read is removed if it has more than 5 N bases
  (positions where the sequencer could not tell which base it was)
\item
  \texttt{-e}: the minimum average quality of a read. The default is 0,
  so this filter is off unless you switch it on
\item
  \texttt{-l}: a read is removed if it is shorter than 15 bases
\end{itemize}

So yes: even without options, fastp filtered the reads on day 2. Your
fastp report from day 2 shows how many reads were removed, and by which
filter.

Filtering is not the only thing that is on by default. The same help
page says that adapter trimming is on by default (\texttt{-A} switches
it off), while quality trimming with \texttt{-\/-cut\_front},
\texttt{-\/-cut\_tail} or \texttt{-\/-cut\_right} is off until you ask
for it. So ``fastp trims and filters my reads'' is only half true, and
which half depends on the option. This decides what ``I ran fastp with
default settings'' actually means, for example in the methods section of
a paper.

\textbf{What two LLMs answered (question 1, October 2026)}

I asked this question to UvA AI Chat (with its default model) and to
ChatGPT, and compared both answers with the help page of fastp 1.4.0. I
used ChatGPT on purpose, without any memory of earlier chats, because it
knows nothing about me. UvA AI Chat does remember my earlier
conversations (see point 5 below). As a student, use UvA AI Chat (see
the UvA rules above).

\begin{enumerate}
\def\labelenumi{\arabic{enumi}.}
\tightlist
\item
  Both got every default value right, but they also didn't give me the
  full output of the help function. Checking this took me a few minutes,
  and that is the point: often the LLM is right and often it might only
  give you a partial answer based on what you asked. A quick look at the
  help page is how you know
\item
  Quality trimming was hidden in one answer, in a side note near the
  end, and missing from the other. Both answered the question I asked,
  which was about filtering. But for most of us, another relevant
  question is ``were my reads trimmed?''. With Illumina data, the
  quality often drops towards the end of the read, so you may want to
  cut those ends off. To ask about that, though, you first need to know
  it. As I mentioned in point 1, an LLM answers only what you ask. It
  may tell you only part of what matters for your data, and nothing in
  the answer tells you that something is missing. The help page doesn't
  depend on your question: it lists every option, including the ones
  that are off and the ones you didn't know to ask about.
\item
  Both wrote that fastp removes reads shorter than 15 bases ``after
  trimming''. That may well be true, but the help page doesn't say it.
  One answer also stayed vague (``fastp \emph{can} discard a pair
  \ldots{} depending on \ldots{}'') about what happens to the other read
  of a pair when one read fails, while the help page says clearly that
  it is discarded by default
\item
  Neither answer said which version of fastp it described
\item
  One answer ended with advice tailored to one of my own research
  projects, which I had not mentioned in the question. It assumed a
  different kind of data than the data I was asking about. The reason
  for that is that depending on your settings an LLM can remember
  earlier conversations. This can be useful for some things but you have
  to decide if this is what you want and whether this fits into the
  privacy rules of your organisation.
\end{enumerate}

\textbf{Differences between versions (question 3)}

Scrolling through two long help pages is tiring. Instead, we can save
both in a file and let the computer compare them. \texttt{diff} compares
two files line by line and prints only the lines that differ:

\begin{Shaded}
\begin{Highlighting}[]
\CommentTok{\# Save the help page of each version in a file}
\ExtensionTok{conda}\NormalTok{ activate fastp\_1.4.0}
\ExtensionTok{fastp} \AttributeTok{{-}{-}help} \OperatorTok{\textgreater{}}\NormalTok{ help\_1.4.0.txt }\DecValTok{2}\OperatorTok{\textgreater{}\&}\DecValTok{1}
\ExtensionTok{conda}\NormalTok{ deactivate}
\ExtensionTok{fastp} \AttributeTok{{-}{-}help} \OperatorTok{\textgreater{}}\NormalTok{ help\_1.0.1.txt }\DecValTok{2}\OperatorTok{\textgreater{}\&}\DecValTok{1}

\CommentTok{\# Print only the lines that differ}
\FunctionTok{diff}\NormalTok{ help\_1.0.1.txt help\_1.4.0.txt}
\end{Highlighting}
\end{Shaded}

On Crunchomics, this prints:

\begin{verbatim}
16c16
<       --stdin                          input from STDIN. If the STDIN is interleaved paired-end FASTQ, please also add --interleaved_in.
---
>       --stdin                          input from STDIN. If the STDIN is interleaved paired-end FASTQ, please also add --interleaved_in. Adapter auto-detection is disabled for STDIN mode
28a29
>       --dimer_max_len                  if the read length is less than or equal to this value after adapter trimming, it is considered an adapter dimer. Requires adapter evidence. (int [=2])
36c37
<       --dup_calc_accuracy              accuracy level to calculate duplication (1~6), higher level uses more memory (1G, 2G, 4G, 8G, 16G, 24G). Default 1 for no-dedup mode, and 3 for dedup mode. (int [=0])
---
>       --dup_calc_accuracy              accuracy level to calculate duplication (1~6), higher level uses more memory (1G, 2G, 4G, 8G, 16G, 32G). Default 1 for no-dedup mode, and 3 for dedup mode. (int [=0])
\end{verbatim}

Lines that start with \texttt{\textless{}} come from the first file
(1.0.1), lines with \texttt{\textgreater{}} from the second file
(1.4.0). The numbers above them say where in each file the lines are.

So for the filtering options, nothing changed between 1.0.1 and 1.4.0.
That is a good outcome too: the point of the check is to \emph{know},
not to find a difference. Three other lines did change. The most
interesting one is \texttt{-\/-dimer\_max\_len}, a new option for
adapter trimming in 1.4.0. Reads of 2 bases or less would be removed by
the length filter anyway, but the help page doesn't say whether the new
option changes anything else. Interestingly, the fastp README on GitHub
for version 1.4.0 doesn't list this option at all (checked in October
2026). So even the official documentation online can lag behind the
program.

\textbf{Key points}

\begin{itemize}
\tightlist
\item
  Often the LLM is right but it might also blur things in its response.
  A quick look at the help page is how you \emph{know}
\item
  An LLM answers only what you ask, and it won't tell you what you
  didn't ask. The help page lists every option, also the ones that are
  off
\item
  ``Default settings'' only means something if you know what is on and
  what is off
\item
  If the LLM, the online documentation and your help page differ, the
  help page of your own installation wins
\end{itemize}
\end{quote}

\end{tcolorbox}

\begin{tcolorbox}[enhanced jigsaw, colframe=quarto-callout-tip-color-frame, opacitybacktitle=0.6, colbacktitle=quarto-callout-tip-color!10!white, title=\textcolor{quarto-callout-tip-color}{\faLightbulb}\hspace{0.5em}{Tip: An LLM that looks something up has not always read the page
(optional)}, rightrule=.15mm, coltitle=black, bottomtitle=1mm, opacityback=0, toptitle=1mm, titlerule=0mm, arc=.35mm, toprule=.15mm, bottomrule=.15mm, left=2mm, colback=white, leftrule=.75mm, breakable]

Some LLM tools can search the web or open a link you give them. That
sounds like they read the page, but often they don't read all of it.
Some tools pass the page through a second, smaller model that summarises
it, and the main model only sees that summary. Claude Code works like
this, for example (checked in October 2026; other tools may work
differently). An answer can also be based only on the short snippets
that a search engine shows for each result.

There is a good reason for this: a long web page is a lot of text to
process, and a summary makes the answer faster and cheaper. So I don't
expect this to disappear, even if the details change from tool to tool.
And for many questions, a summary is all you need, for example ``What is
fastp used for?'' or ``Which tools are often used to trim reads?''.

It becomes a problem when the detail is the point. A summary can drop a
detail, mix up two options, or smooth over exactly the sentence that
matters, and it reads just as convincingly as a correct answer. So when
you need a default value, an option name or the version a page
describes, open the page or the help page yourself and read the lines in
question. Ask the LLM to quote the exact sentence it relies on, so you
can find it and check it.

\end{tcolorbox}

\section{Using an LLM to understand
code}\label{using-an-llm-to-understand-code}

This is where I find LLMs most useful. You can paste a command and ask
what each part does. This works well with free chat tools.

Some questions that work well:

\begin{verbatim}
Explain this bash command line by line. I am a biologist and new to the command line:

<paste the command here>
\end{verbatim}

\begin{verbatim}
What happens to each file when I run this loop?
What happens if one of the input files is missing?
\end{verbatim}

\begin{verbatim}
What is the difference between the -o and the -h option of fastp?
\end{verbatim}

Keep in mind that the answer you get can still be wrong. An explanation
is easier to check than a full script, though: you can test each claim
with a small command or by looking at the help function yourself.

\begin{tcolorbox}[enhanced jigsaw, colframe=quarto-callout-color-frame, opacitybacktitle=0.6, colbacktitle=quarto-callout-color!10!white, title=\textcolor{quarto-callout-color}{\faInfo}\hspace{0.5em}{Exercise}, rightrule=.15mm, coltitle=black, bottomtitle=1mm, opacityback=0, toptitle=1mm, titlerule=0mm, arc=.35mm, toprule=.15mm, bottomrule=.15mm, left=2mm, colback=white, leftrule=.75mm, breakable]

In the optional box at the end of the fastp section on day 2, we used
this awk command (if you skipped that box, you don't have
\texttt{stats\_trimmed.txt}, but you can still do the checks in the
answer on \texttt{stats\_raw.txt}):

\begin{Shaded}
\begin{Highlighting}[]
\FunctionTok{awk} \AttributeTok{{-}F}\StringTok{\textquotesingle{}\textbackslash{}t\textquotesingle{}} \StringTok{\textquotesingle{}FNR \textgreater{} 1 \{print $1, $4, $7\}\textquotesingle{}}\NormalTok{ results/seqkit/stats\_raw.txt results/seqkit/stats\_trimmed.txt}
\end{Highlighting}
\end{Shaded}

\begin{enumerate}
\def\labelenumi{\arabic{enumi}.}
\tightlist
\item
  Paste the command into UvA AI Chat (or another chat tool that you are
  allowed to use) and ask it to explain each part. This exercise is
  about practicing checking an answer from an LLM, not about getting the
  code answer.
\item
  Pick one thing the explanation claims, for example what
  \texttt{FNR\ \textgreater{}\ 1} does. How can you check this claim
  yourself, with the commands from day 1?
\end{enumerate}

\begin{quote}
\textbf{Click me to see an answer}

For question 2: the explanation should say that
\texttt{FNR\ \textgreater{}\ 1} skips the first line of each file, the
header. You can check this by comparing the input with the output:

\begin{Shaded}
\begin{Highlighting}[]
\CommentTok{\# Look at the first lines of the input: the first line is the header}
\FunctionTok{head} \AttributeTok{{-}n}\NormalTok{ 3 results/seqkit/stats\_raw.txt}

\CommentTok{\# Run awk on one file only: the header line should be missing}
\FunctionTok{awk} \AttributeTok{{-}F}\StringTok{\textquotesingle{}\textbackslash{}t\textquotesingle{}} \StringTok{\textquotesingle{}FNR \textgreater{} 1 \{print $1, $4, $7\}\textquotesingle{}}\NormalTok{ results/seqkit/stats\_raw.txt }\KeywordTok{|} \FunctionTok{head} \AttributeTok{{-}n}\NormalTok{ 3}

\CommentTok{\# Count the lines: the awk output should have one line less than the input}
\FunctionTok{wc} \AttributeTok{{-}l}\NormalTok{ results/seqkit/stats\_raw.txt}
\FunctionTok{awk} \AttributeTok{{-}F}\StringTok{\textquotesingle{}\textbackslash{}t\textquotesingle{}} \StringTok{\textquotesingle{}FNR \textgreater{} 1 \{print $1, $4, $7\}\textquotesingle{}}\NormalTok{ results/seqkit/stats\_raw.txt }\KeywordTok{|} \FunctionTok{wc} \AttributeTok{{-}l}
\end{Highlighting}
\end{Shaded}

You don't need to know awk to do this. You only need to know what the
output \emph{should} look like. This is also a good example for how
knowing basic bash commands allows you to perform the suggested checks
yourself, no LLM needed.

Another claim worth checking is what columns 1, 4 and 7 contain. When I
asked ChatGPT (October 2026), its first answer showed an example row
where column 7 was \texttt{Q30}. Further down, it did say that its
examples were only illustrative. But if you skim a long answer, it is
easy to miss that sentence and remember the example as a fact.

The header line of the file tells you what each column really is:

\begin{Shaded}
\begin{Highlighting}[]
\CommentTok{\# Print only the header line}
\FunctionTok{head} \AttributeTok{{-}n}\NormalTok{ 1 results/seqkit/stats\_raw.txt}
\end{Highlighting}
\end{Shaded}

This prints:

\begin{verbatim}
file    format  type    num_seqs        sum_len min_len avg_len max_len Q1      Q2      Q3      sum_gap N50       Q20(%)  Q30(%)  AvgQual GC(%)
\end{verbatim}

If you count the columns, column 1 is \texttt{file}, column 4 is
\texttt{num\_seqs} and column 7 is \texttt{avg\_len}, the average read
length. So the command prints the file name, the number of reads and the
average read length, not a quality value.

What makes this mistake tricky is that the answer started out right. The
command mentions \texttt{results/seqkit/}, and the example file in the
answer began with the real seqkit column names: \texttt{file},
\texttt{format}, \texttt{type}, \texttt{num\_seqs}. The example values
for columns 2 to 6 also fit those columns. Only column 7 was suddenly a
quality value. When the first part of an answer checks out, it is
tempting to trust the rest as well. But an LLM doesn't get more reliable
further into an answer, and one correct detail says nothing about the
next one.

So check the detail that matters to you, not the first one you can
check. And whenever an answer depends on your own files, look at the
file yourself, for example with \texttt{head}.
\end{quote}

\end{tcolorbox}

\begin{tcolorbox}[enhanced jigsaw, colframe=quarto-callout-tip-color-frame, opacitybacktitle=0.6, colbacktitle=quarto-callout-tip-color!10!white, title=\textcolor{quarto-callout-tip-color}{\faLightbulb}\hspace{0.5em}{Tip: When the answer is longer than you want to read (optional)}, rightrule=.15mm, coltitle=black, bottomtitle=1mm, opacityback=0, toptitle=1mm, titlerule=0mm, arc=.35mm, toprule=.15mm, bottomrule=.15mm, left=2mm, colback=white, leftrule=.75mm, breakable]

LLM answers are often long. When I asked ChatGPT (October 2026) to
explain the one-line awk command from the exercise above line by line, I
got seven sections, two tables, example files and a summary at the end.
I used ChatGPT without any memory of earlier chats on purpose, so it
knew nothing about me except what I wrote in the question. Most of it
was correct (the answer above shows one detail that was not). But the
part I actually needed fits in one sentence, and it was near the top:

\begin{verbatim}
Read two files, skip the first line of each file, and print columns 1, 4, and 7 from every remaining line.
\end{verbatim}

The rest explained things like what the folder \texttt{results/} is, and
why a biologist might want to use this command. The few details that
really matter, such as how \texttt{FNR} is different from \texttt{NR},
were somewhere in the middle. If you are new to the command line, it is
hard to tell which parts are important and which are filler. And
scrolling through a long, well-written answer can feel like learning,
even when not much of it sticks.

You don't need a perfect question to avoid this. You can simply say what
is wrong with the answer. I replied that most students would not read
through this much text, and the second answer was much shorter. Its core
was this table:

\begin{verbatim}
Part                Meaning
awk                 Reads and processes text, one line at a time.
-F'\t'              Treats tabs as column separators.
FNR > 1             Skips the first line of each file (usually the header).
print $1, $4, $7    Prints columns 1, 4, and 7.
stats_raw.txt       First input file.
stats_trimmed.txt   Second input file.
\end{verbatim}

That fits on one screen, and every line is something you can check.
Other short follow-ups that work:

\begin{verbatim}
Shorter please, max 5 lines.
\end{verbatim}

\begin{verbatim}
Only explain what FNR > 1 does.
\end{verbatim}

\begin{verbatim}
Which part of your answer should I check myself, and how?
\end{verbatim}

Asking about one part at a time also makes each answer easier to check.

\end{tcolorbox}

\section{Checking code with the commands from day
1}\label{checking-code-with-the-commands-from-day-1}

This is the main point of this page. Not all code needs the same level
of understanding:

\begin{itemize}
\tightlist
\item
  \textbf{The analysis itself: understand every option.} The options of
  a tool like fastp decide what happens to your data: which reads are
  removed, and how much of each read is trimmed. These are scientific
  choices. You need to know what each option does and why you use it,
  whether you or an LLM wrote the command. Look them up with
  \texttt{fastp\ -\/-help} or in the manual, and ask an LLM to explain
  an option if the manual is unclear
\item
  \textbf{The code around it: know what it should do, then check.} Some
  code only organises the work: an \texttt{if}/\texttt{else} block that
  skips finished samples, or a \texttt{sed} or \texttt{awk} command that
  picks lines from a file. You may not understand every detail of how it
  works. But you must know what it is supposed to do to your files, and
  check that it really does
\end{itemize}

For these checks, the commands from day 1 are enough:

\begin{longtable}[]{@{}
  >{\raggedright\arraybackslash}p{(\linewidth - 2\tabcolsep) * \real{0.4583}}
  >{\raggedright\arraybackslash}p{(\linewidth - 2\tabcolsep) * \real{0.5417}}@{}}
\toprule\noalign{}
\begin{minipage}[b]{\linewidth}\raggedright
Question
\end{minipage} & \begin{minipage}[b]{\linewidth}\raggedright
Command
\end{minipage} \\
\midrule\noalign{}
\endhead
\bottomrule\noalign{}
\endlastfoot
Did I get all the output files I expected? &
\texttt{ls\ -l\ results/fastp},
\texttt{ls\ results/fastp\ \textbackslash{}\textbar{}\ wc\ -l} \\
Is any file empty? & \texttt{ls\ -l} (look for a size of 0) \\
Which samples will the loop work on? & \texttt{cat\ samples.txt},
\texttt{wc\ -l\ samples.txt} \\
What does a wildcard match? &
\texttt{ls\ data/seq\_project/*/*\_R1.fastq.gz} \\
How many reads are in a file? &
\texttt{zcat\ file.fastq.gz\ \textbackslash{}\textbar{}\ wc\ -l}
(divided by 4) \\
Does the output look right? & \texttt{head}, \texttt{less\ -S} \\
Did every sample show up in the log? &
\texttt{grep\ -c\ "Start\ fastp"\ logs/fastp\_*.out} \\
What will the loop do, before it does it? & put \texttt{echo} in front
of the command (a dry run) \\
\end{longtable}

Before you run anything, think about what you expect: how many samples,
how many files, roughly how many reads. Then compare. If the numbers
don't match, you have found a problem, even if you can't yet say where
it is.

\begin{tcolorbox}[enhanced jigsaw, colframe=quarto-callout-color-frame, opacitybacktitle=0.6, colbacktitle=quarto-callout-color!10!white, title=\textcolor{quarto-callout-color}{\faInfo}\hspace{0.5em}{Exercise}, rightrule=.15mm, coltitle=black, bottomtitle=1mm, opacityback=0, toptitle=1mm, titlerule=0mm, arc=.35mm, toprule=.15mm, bottomrule=.15mm, left=2mm, colback=white, leftrule=.75mm, breakable]

Imagine you ask an LLM how to count the reads in a fastq file, and it
suggests:

\begin{Shaded}
\begin{Highlighting}[]
\FunctionTok{zcat}\NormalTok{ data/seq\_project/barcode001\_User1\_ITS\_1\_L001/Sample{-}DUMMY1\_R1.fastq.gz }\KeywordTok{|} \FunctionTok{grep} \AttributeTok{{-}c} \StringTok{"\^{}@"}
\end{Highlighting}
\end{Shaded}

The idea: each read starts with a header line that begins with
\texttt{@}, so counting those lines should count the reads
(\texttt{\^{}@} means ``a line that starts with @'').

\begin{enumerate}
\def\labelenumi{\arabic{enumi}.}
\tightlist
\item
  Run the command. How many reads does it report?
\item
  Count the reads in the same file the way we did on day 1. Do the two
  numbers match?
\item
  Can you think of a reason why they don't match? Hint: look at the
  fastq format again. Which other line of a read could start with
  \texttt{@}?
\end{enumerate}

\begin{quote}
\textbf{Click me to see an answer}

On my computer (WSL on Windows), the \texttt{grep} command prints:

\begin{verbatim}
44882
\end{verbatim}

The day 1 way:

\begin{Shaded}
\begin{Highlighting}[]
\FunctionTok{zcat}\NormalTok{ data/seq\_project/barcode001\_User1\_ITS\_1\_L001/Sample{-}DUMMY1\_R1.fastq.gz }\KeywordTok{|} \FunctionTok{wc} \AttributeTok{{-}l}
\end{Highlighting}
\end{Shaded}

This prints:

\begin{verbatim}
179384
\end{verbatim}

179,384 lines divided by 4 is 44,846 reads. So the \texttt{grep} command
counts 36 reads too many.

The reason: the fourth line of each read holds the quality scores, and
\texttt{@} is one of the characters used for a quality score. In 36
reads, the quality line happens to start with \texttt{@}, so
\texttt{grep} counts these lines as well.

The command looked sensible, and the number looked plausible. Only the
comparison with a second, simple count showed the problem.
\end{quote}

\end{tcolorbox}

\section{Worked example: an LLM writes a fastp job
script}\label{worked-example-an-llm-writes-a-fastp-job-script}

Let's look at a realistic example. I asked an LLM:

\begin{verbatim}
I work with samples S1-S100. I want to run fastp for quality trimming and
remove reads below an average quality score of 20 and use the same cutoff
for quality trimming of the reads themselves. Can you write a script to
process all these samples at once on a HPC with slurm?
\end{verbatim}

The prompt could be better, but it is roughly what many of us would
type. This is the script it wrote (I only shortened the decorative lines
in the comments):

\begin{Shaded}
\begin{Highlighting}[]
\CommentTok{\#!/bin/bash}
\CommentTok{\#SBATCH {-}{-}job{-}name=fastp\_trim}
\CommentTok{\#SBATCH {-}{-}array=1{-}100}
\CommentTok{\#SBATCH {-}{-}cpus{-}per{-}task=8}
\CommentTok{\#SBATCH {-}{-}mem=16G}
\CommentTok{\#SBATCH {-}{-}time=04:00:00}
\CommentTok{\#SBATCH {-}{-}output=logs/fastp\_\%A\_\%a.out}
\CommentTok{\#SBATCH {-}{-}error=logs/fastp\_\%A\_\%a.err}

\CommentTok{\# {-}{-} Paths {-}{-}}
\VariableTok{DATA\_DIR}\OperatorTok{=}\StringTok{"data"}
\VariableTok{OUT\_DIR}\OperatorTok{=}\StringTok{"trimmed"}
\FunctionTok{mkdir} \AttributeTok{{-}p} \StringTok{"}\VariableTok{$\{OUT\_DIR\}}\StringTok{"}\NormalTok{ logs}

\CommentTok{\# {-}{-} Sample name derived from array index {-}{-}}
\VariableTok{SAMPLE}\OperatorTok{=}\StringTok{"S}\VariableTok{$\{SLURM\_ARRAY\_TASK\_ID\}}\StringTok{"}

\VariableTok{R1}\OperatorTok{=}\StringTok{"}\VariableTok{$\{DATA\_DIR\}}\StringTok{/}\VariableTok{$\{SAMPLE\}}\StringTok{\_R1.fastq.gz"}
\VariableTok{R2}\OperatorTok{=}\StringTok{"}\VariableTok{$\{DATA\_DIR\}}\StringTok{/}\VariableTok{$\{SAMPLE\}}\StringTok{\_R2.fastq.gz"}

\CommentTok{\# {-}{-} Guard: skip if input files are missing {-}{-}}
\ControlFlowTok{if} \KeywordTok{[[} \OtherTok{!} \OtherTok{{-}f} \StringTok{"}\VariableTok{$\{R1\}}\StringTok{"} \KeywordTok{||} \OtherTok{!} \OtherTok{{-}f} \StringTok{"}\VariableTok{$\{R2\}}\StringTok{"} \KeywordTok{]];} \ControlFlowTok{then}
    \BuiltInTok{echo} \StringTok{"ERROR: Input files not found for }\VariableTok{$\{SAMPLE\}}\StringTok{: }\VariableTok{$\{R1\}}\StringTok{ / }\VariableTok{$\{R2\}}\StringTok{"} \OperatorTok{\textgreater{}\&}\DecValTok{2}
    \BuiltInTok{exit}\NormalTok{ 1}
\ControlFlowTok{fi}

\CommentTok{\# {-}{-} Run fastp {-}{-}}
\ExtensionTok{fastp} \DataTypeTok{\textbackslash{}}
    \AttributeTok{{-}{-}in1}  \StringTok{"}\VariableTok{$\{R1\}}\StringTok{"} \DataTypeTok{\textbackslash{}}
    \AttributeTok{{-}{-}in2}  \StringTok{"}\VariableTok{$\{R2\}}\StringTok{"} \DataTypeTok{\textbackslash{}}
    \AttributeTok{{-}{-}out1} \StringTok{"}\VariableTok{$\{OUT\_DIR\}}\StringTok{/}\VariableTok{$\{SAMPLE\}}\StringTok{\_R1\_trimmed.fastq.gz"} \DataTypeTok{\textbackslash{}}
    \AttributeTok{{-}{-}out2} \StringTok{"}\VariableTok{$\{OUT\_DIR\}}\StringTok{/}\VariableTok{$\{SAMPLE\}}\StringTok{\_R2\_trimmed.fastq.gz"} \DataTypeTok{\textbackslash{}}
    \AttributeTok{{-}{-}average\_qual}\NormalTok{        20 }\DataTypeTok{\textbackslash{} }  \CommentTok{\# discard reads with mean quality \textless{} 20}
    \ExtensionTok{{-}{-}cut\_mean\_quality}\NormalTok{    20 }\DataTypeTok{\textbackslash{} }  \CommentTok{\# sliding{-}window trimming threshold}
    \ExtensionTok{{-}{-}cut\_front}              \DataTypeTok{\textbackslash{} }  \CommentTok{\# trim low{-}quality bases from 5\textquotesingle{} end}
    \ExtensionTok{{-}{-}cut\_tail}               \DataTypeTok{\textbackslash{} }  \CommentTok{\# trim low{-}quality bases from 3\textquotesingle{} end}
    \ExtensionTok{{-}{-}thread}\NormalTok{              8  }\DataTypeTok{\textbackslash{}}
    \AttributeTok{{-}{-}json} \StringTok{"}\VariableTok{$\{OUT\_DIR\}}\StringTok{/}\VariableTok{$\{SAMPLE\}}\StringTok{\_fastp.json"} \DataTypeTok{\textbackslash{}}
    \AttributeTok{{-}{-}html} \StringTok{"}\VariableTok{$\{OUT\_DIR\}}\StringTok{/}\VariableTok{$\{SAMPLE\}}\StringTok{\_fastp.html"}

\BuiltInTok{echo} \StringTok{"Done: }\VariableTok{$\{SAMPLE\}}\StringTok{"}
\end{Highlighting}
\end{Shaded}

It looks professional: tidy comments, a check on the input files,
separate log files per sample. It uses a \textbf{job array}, which we
only saw in an optional box on day 2: SLURM starts the same script many
times (here 100 times), and each copy gets its own number in the
variable \texttt{SLURM\_ARRAY\_TASK\_ID}. The script uses the long names
of the fastp options (\texttt{-\/-in1} instead of \texttt{-i},
\texttt{-\/-out1} instead of \texttt{-o}); both forms do the same thing.

\begin{tcolorbox}[enhanced jigsaw, colframe=quarto-callout-color-frame, opacitybacktitle=0.6, colbacktitle=quarto-callout-color!10!white, title=\textcolor{quarto-callout-color}{\faInfo}\hspace{0.5em}{Exercise}, rightrule=.15mm, coltitle=black, bottomtitle=1mm, opacityback=0, toptitle=1mm, titlerule=0mm, arc=.35mm, toprule=.15mm, bottomrule=.15mm, left=2mm, colback=white, leftrule=.75mm, breakable]

Read the script carefully, as if you were going to run it on your own
data. Write down everything you would check or change. Some questions to
guide you:

\begin{enumerate}
\def\labelenumi{\arabic{enumi}.}
\tightlist
\item
  Which fastp options will fastp actually receive? Look closely at the
  lines with a comment at the end
\item
  What does the script assume about your file names and folders? Would
  it work with the files from this workshop?
\item
  Which fastp is used when the job runs on Crunchomics?
\item
  Where do the numbers for CPUs, memory and time come from?
\item
  Compare it with the fastp script from day 2. Is it wrong, or just
  different?
\end{enumerate}

\begin{quote}
\textbf{Click me to see an answer}

I sort the problems from most to least serious.

\textbf{1. It runs, but does the wrong thing}

\begin{itemize}
\tightlist
\item
  \textbf{The comments at the end of the lines break the command.} The
  \texttt{\textbackslash{}} only joins two lines when it is the very
  last character on the line (we saw this on day 2). Here, spaces and a
  comment follow it. The \texttt{\textbackslash{}} then only protects
  the space after it, and the command ends at the end of that line.
  fastp receives the input and output files,
  \texttt{-\/-average\_qual\ 20}, and one extra argument that is just a
  space. Bash then tries to run the lines that follow as commands of
  their own, and each attempt fails with \texttt{command\ not\ found}.
  The \texttt{-\/-thread} line has no comment, so its
  \texttt{\textbackslash{}} works again, and \texttt{-\/-thread},
  \texttt{-\/-json} and \texttt{-\/-html} end up together in one failing
  command. So fastp never gets the trimming options,
  \texttt{-\/-thread}, \texttt{-\/-json} or \texttt{-\/-html}. I missed
  this myself when I first read the script, which shows that reading is
  not enough: you also have to run it on one sample and look at the
  output. See the next exercise for a small test.
\item
  \textbf{The job reports success anyway.} The last line,
  \texttt{echo\ "Done"}, always works. A job counts as successful when
  its last command works, so SLURM reports every task as finished
  without problems. The errors are only visible in the \texttt{.err}
  files. (With \texttt{set\ -euo\ pipefail} at the top of the script, it
  would have stopped at the first \texttt{command\ not\ found}, and
  SLURM would have reported the task as failed. fastp would still have
  run once with the incomplete options before that. See the tip on
  \texttt{set\ -euo\ pipefail} further down.)
\item
  \textbf{The \texttt{logs} folder is made too late.} SLURM opens the
  log files before the script starts, so the \texttt{mkdir} for
  \texttt{logs} inside the script comes too late to matter. On
  Crunchomics, SLURM then makes a missing \texttt{logs} folder itself,
  but inside the folder where you typed \texttt{sbatch}, which is not
  always where you expect your logs (we saw this on day 2). Make the
  folder before you run \texttt{sbatch}, as we did on day 2, so that you
  know where your logs end up.
\end{itemize}

\textbf{2. It makes assumptions about your data}

\begin{itemize}
\tightlist
\item
  \textbf{The file names are invented.}
  \texttt{S\$\{SLURM\_ARRAY\_TASK\_ID\}} only works because the prompt
  said ``S1-S100''. Real file names rarely look like that. In this
  workshop, the files are called \texttt{Sample-DUMMY1\_R1.fastq.gz} and
  are in separate folders. A list file like \texttt{samples.txt} works
  with any names, and you can open it to see which samples will be
  processed.
\item
  \textbf{It chose a trimming method for you.} The prompt was not clear
  about how to trim, so the LLM chose trimming from both ends of the
  read (without asking me for clarification). Is that what you want?
  fastp 1.4.0 offers three methods, each with a window of 4 bases that
  moves along the read: \texttt{-\/-cut\_front} (from the start of the
  read, removes bases until the quality is good enough),
  \texttt{-\/-cut\_tail} (the same, from the end of the read) and
  \texttt{-\/-cut\_right} (from the start; at the first window with low
  quality, it removes that window and everything after it). Which one
  fits depends on your data. \texttt{-\/-cut\_mean\_quality} only sets
  the quality threshold that these three methods use. On its own, it
  does nothing: fastp only trims by quality if you also add
  \texttt{-\/-cut\_front}, \texttt{-\/-cut\_tail} or
  \texttt{-\/-cut\_right}. (And 20 is already the default threshold, so
  the script would trim the same way without this line.) And fastp
  already filters reads by default, also without any of these options
  (see the exercise on default filtering above). Check
  \texttt{fastp\ -\/-help} for the version you use, and ask yourself (or
  the LLM) why these options and not others.
\item
  \textbf{The resources are guesses.} 8 CPUs, 16 GB of memory and 4
  hours per sample are not based on your data. Asking for too much time
  means a longer wait in the queue. Run one sample first, then check
  what it actually used (the optional \texttt{sacct} box on day 2).
\item
  \textbf{No software is activated.} There is no
  \texttt{conda\ activate}, so the script doesn't say which fastp it
  needs. Which one the job uses then depends on your terminal at the
  moment you submit the job: SLURM normally copies its settings,
  including the PATH, into the job. If the \texttt{fastp\_1.4.0}
  environment happened to be active, the job uses fastp 1.4.0. If not,
  it uses the fastp 1.0.1 of Crunchomics. Your results should not depend
  on such an accident.
\end{itemize}

\textbf{3. It is harder to read than needed}

\begin{itemize}
\tightlist
\item
  \texttt{DATA\_DIR} and \texttt{OUT\_DIR} hide two short folder names.
  You have to look them up to know where the files go. Variables like
  \texttt{SAMPLE}, \texttt{R1} and \texttt{R2}, however, are useful,
  because they are used many times.
\item
  The one place where a variable \emph{would} help is the number 8: it
  appears twice (\texttt{-\/-cpus-per-task=8} and
  \texttt{-\/-thread\ 8}). If you change one, you have to remember the
  other. \texttt{-\/-thread\ \$\{SLURM\_CPUS\_PER\_TASK\}} would always
  use the number of CPUs you requested.
\item
  The check on the input files looks safe, but it is easy to overlook
  when it fires. If a file is missing, only that one task stops. Its
  error message ends up in one \texttt{.err} file among many, and the
  other tasks finish normally. With 4 samples you might notice. With
  1000 samples, one failed task is easy to miss, and you continue with
  999 samples without knowing it. I personally prefer checks that fail
  loudly, before anything runs: check that all input files exist
  \emph{before} you submit the job, with \texttt{wc\ -l\ samples.txt}
  and a loop that runs \texttt{ls} on every R1 and R2 file. If a file is
  missing, you see it right away and the job never starts. And after the
  job, count the output files, whatever checks the script contains.
\end{itemize}

\textbf{4. Different from day 2: wrong, or just different?}

The LLM used a job array, while we used a for-loop inside one job
script. Both are fine. An array runs the samples at the same time, a
loop runs them one after the other. Something different from what you
learned is not wrong by itself. But you should understand it before you
use it. Also be careful with arrays on a shared HPC:
\texttt{-\/-array=1-100} with 8 CPUs each can ask for 800 CPUs at once.
\texttt{-\/-array=1-100\%8} runs at most 8 tasks at the same time.
\end{quote}

\end{tcolorbox}

You can test the first problem without fastp and without an HPC. Put
\texttt{echo} in front of the command, as in a dry run, and add comments
at the end of the lines like the LLM did. Save these lines in a script
called \texttt{comment\_test.sh}:

\begin{Shaded}
\begin{Highlighting}[]
\BuiltInTok{echo}\NormalTok{ fastp }\DataTypeTok{\textbackslash{}}
    \AttributeTok{{-}{-}in1}\NormalTok{ Sample{-}DUMMY1\_R1.fastq.gz }\DataTypeTok{\textbackslash{} }  \CommentTok{\# the R1 file}
    \ExtensionTok{{-}{-}in2}\NormalTok{ Sample{-}DUMMY1\_R2.fastq.gz }\DataTypeTok{\textbackslash{} }  \CommentTok{\# the R2 file}
    \ExtensionTok{{-}{-}thread}\NormalTok{ 2}
\BuiltInTok{echo} \StringTok{"Done"}
\end{Highlighting}
\end{Shaded}

Then run it with \texttt{bash\ comment\_test.sh}. On my computer (WSL
with Ubuntu), this prints:

\begin{verbatim}
fastp --in1 Sample-DUMMY1_R1.fastq.gz
comment_test.sh: line 3: --in2: command not found
comment_test.sh: line 4: --thread: command not found
Done
\end{verbatim}

The first line shows the full command that fastp would receive: only the
R1 file. The R2 file and \texttt{-\/-thread} never reach it, and bash
simply continues with the next lines. In the worked example, this is
what makes the mistake dangerous: fastp still gets the input and output
files and \texttt{-\/-average\_qual\ 20}, so it runs, writes trimmed
files, and SLURM reports success. The output looks normal, but it was
made without the trimming options you asked for.

\section{Less code is often better}\label{less-code-is-often-better}

LLMs have learned from a lot of code written for large, automated
software systems. There, a script should keep running no matter what
goes wrong. In a scientific analysis, you usually want the opposite: if
something goes wrong, you want to know.

A common example is a check that skips a sample when its input file is
missing. Here, I made a typo in \texttt{samples.txt}
(\texttt{Sample-DUMY3} instead of \texttt{Sample-DUMMY3}) and ran a loop
in the style that LLMs often write. To keep it short, the loop only
prints which sample it would work on:

\begin{Shaded}
\begin{Highlighting}[]
\ControlFlowTok{for}\NormalTok{ sample }\KeywordTok{in} \VariableTok{$(}\FunctionTok{cat}\NormalTok{ samples.txt}\VariableTok{)}\KeywordTok{;} \ControlFlowTok{do}
    \VariableTok{R1}\OperatorTok{=}\VariableTok{$(}\FunctionTok{ls}\NormalTok{ data/seq\_project/}\PreprocessorTok{*}\NormalTok{/}\VariableTok{$\{sample\}}\NormalTok{\_R1.fastq.gz }\DecValTok{2}\OperatorTok{\textgreater{}}\NormalTok{/dev/null}\VariableTok{)}
    \ControlFlowTok{if} \BuiltInTok{[} \OtherTok{!} \OtherTok{{-}f} \StringTok{"}\VariableTok{$R1}\StringTok{"} \BuiltInTok{]}\KeywordTok{;} \ControlFlowTok{then}
        \ControlFlowTok{continue}
    \ControlFlowTok{fi}
    \BuiltInTok{echo} \StringTok{"Processing }\VariableTok{$sample}\StringTok{"}
\ControlFlowTok{done}
\BuiltInTok{echo} \StringTok{"All samples done"}
\end{Highlighting}
\end{Shaded}

On my computer, this prints:

\begin{verbatim}
Processing Sample-DUMMY1
Processing Sample-DUMMY2
Processing Sample-DUMMY4
All samples done
\end{verbatim}

\texttt{continue} tells the loop to skip the rest of this round and go
on with the next sample. The loop skipped sample 3 without a single
word, and then reported that all samples were done. If this were a long
fastp job, the results would look complete, but cover only 3 of the 4
samples. \texttt{2\textgreater{}/dev/null} makes it worse: it throws
away the error message of \texttt{ls}.

Compare this with the simple loop from day 2 (here with \texttt{ls}
instead of fastp):

\begin{Shaded}
\begin{Highlighting}[]
\ControlFlowTok{for}\NormalTok{ sample }\KeywordTok{in} \VariableTok{$(}\FunctionTok{cat}\NormalTok{ samples.txt}\VariableTok{)}\KeywordTok{;} \ControlFlowTok{do}
    \BuiltInTok{echo} \StringTok{"Start: }\VariableTok{$sample}\StringTok{"}
    \FunctionTok{ls}\NormalTok{ data/seq\_project/}\PreprocessorTok{*}\NormalTok{/}\VariableTok{$\{sample\}}\NormalTok{\_R1.fastq.gz}
\ControlFlowTok{done}
\end{Highlighting}
\end{Shaded}

This prints:

\begin{verbatim}
Start: Sample-DUMMY1
data/seq_project/barcode001_User1_ITS_1_L001/Sample-DUMMY1_R1.fastq.gz
Start: Sample-DUMMY2
data/seq_project/barcode002_User1_ITS_9_L001/Sample-DUMMY2_R1.fastq.gz
Start: Sample-DUMY3
ls: cannot access 'data/seq_project/*/Sample-DUMY3_R1.fastq.gz': No such file or directory
Start: Sample-DUMMY4
data/seq_project/barcode004_User2_ITS_9_L001/Sample-DUMMY4_R1.fastq.gz
\end{verbatim}

The simple loop has no check at all, and that is exactly why the problem
is visible: the error message appears right below the sample name, and
ends up in the log file of a job. Together with a check of the output (4
samples in, 16 files out), you would notice the missing sample right
away.

So knowing your own data and checking it with basic commands often works
better than adding checks to the code. Before you submit a job, check
the inputs with \texttt{cat\ samples.txt}, \texttt{wc\ -l\ samples.txt}
and an \texttt{ls} loop like the one above. After the job, count the
outputs.

\begin{tcolorbox}[enhanced jigsaw, colframe=quarto-callout-tip-color-frame, opacitybacktitle=0.6, colbacktitle=quarto-callout-tip-color!10!white, title=\textcolor{quarto-callout-tip-color}{\faLightbulb}\hspace{0.5em}{Tip: Skipping samples that are already done (optional)}, rightrule=.15mm, coltitle=black, bottomtitle=1mm, opacityback=0, toptitle=1mm, titlerule=0mm, arc=.35mm, toprule=.15mm, bottomrule=.15mm, left=2mm, colback=white, leftrule=.75mm, breakable]

Sometimes a check in the loop does make sense. Say a job stopped
halfway, and you don't want to run fastp again for the samples that are
already done. Then write the check so that you can read it, and so that
every skip shows up in the log:

\begin{Shaded}
\begin{Highlighting}[]
\ControlFlowTok{for}\NormalTok{ sample }\KeywordTok{in} \VariableTok{$(}\FunctionTok{cat}\NormalTok{ samples.txt}\VariableTok{)}\KeywordTok{;} \ControlFlowTok{do}
    \ControlFlowTok{if} \BuiltInTok{[} \OtherTok{{-}s}\NormalTok{ results/fastp/}\VariableTok{$\{sample\}}\NormalTok{\_R1.trimmed.fastq.gz }\BuiltInTok{]}\KeywordTok{;} \ControlFlowTok{then}
        \BuiltInTok{echo} \StringTok{"Skip }\VariableTok{$sample}\StringTok{: output already exists"}
    \ControlFlowTok{else}
        \BuiltInTok{echo} \StringTok{"Run fastp for }\VariableTok{$sample}\StringTok{"}
    \ControlFlowTok{fi}
\ControlFlowTok{done}
\end{Highlighting}
\end{Shaded}

\begin{itemize}
\tightlist
\item
  \texttt{if\ {[}\ ...\ {]};\ then\ ...\ else\ ...\ fi}: if the test
  between the square brackets is true, bash runs the commands after
  \texttt{then}. Otherwise, it runs the commands after \texttt{else}.
  \texttt{fi} (if spelled backwards) closes the block
\item
  \texttt{-s}: true if the file exists \emph{and} is not empty. A job
  that crashed can leave an empty file behind, which \texttt{-f} (true
  if the file exists) would accept
\item
  Neither \texttt{-s} nor \texttt{-f} checks that a file is complete. A
  job that ran out of time can leave a half-written file. When in doubt,
  delete the outputs of the sample that was running when the job
  stopped, and run it again
\end{itemize}

Replace the \texttt{echo} after \texttt{else} with the fastp command
once the dry run looks right.

The IBED bioinformatics pages explain this in more detail, each with a
note on code written by LLMs:

\begin{itemize}
\tightlist
\item
  \href{https://scienceparkstudygroup.github.io/ibed-bioinformatics-page/source/core_tools/bash-input-validation.html}{if/else,
  \texttt{continue} and \texttt{break}}: checking input files
\item
  \href{https://scienceparkstudygroup.github.io/ibed-bioinformatics-page/source/core_tools/bash-pipefail.html}{\texttt{set\ -euo\ pipefail}}:
  a line that LLMs often put at the top of a script. It makes a script
  stop at the first command that fails. That sounds like exactly the
  loud failure I recommend on this page, and in the worked example it
  would have helped. But it also stops a script for things that are not
  errors. For example, \texttt{grep} counts as failed when it finds no
  match, even if ``no match'' is a perfectly fine answer. While you
  develop a script, this makes it harder to see what is going on. So
  leave it out while you develop, and only consider adding it once the
  script works and you know which of your commands can ``fail'' without
  a real problem
\end{itemize}

\end{tcolorbox}

\begin{tcolorbox}[enhanced jigsaw, colframe=quarto-callout-color-frame, opacitybacktitle=0.6, colbacktitle=quarto-callout-color!10!white, title=\textcolor{quarto-callout-color}{\faInfo}\hspace{0.5em}{Exercise}, rightrule=.15mm, coltitle=black, bottomtitle=1mm, opacityback=0, toptitle=1mm, titlerule=0mm, arc=.35mm, toprule=.15mm, bottomrule=.15mm, left=2mm, colback=white, leftrule=.75mm, breakable]

Take the LLM script from the worked example, and change it into a script
that works for the data of this workshop, in the style of day 2. Keep
the job array if you like, or use a for-loop.

Then answer: which lines did you remove, and why? Which lines did you
add?

\begin{quote}
\textbf{Click me to see an answer}

There are several good answers. This one keeps the array, and takes the
sample name from \texttt{samples.txt}. Make the \texttt{logs} folder
before you submit it.

\begin{Shaded}
\begin{Highlighting}[]
\CommentTok{\#!/bin/bash}
\CommentTok{\#SBATCH {-}{-}job{-}name=fastp\_array}
\CommentTok{\#SBATCH {-}{-}output=logs/fastp\_array\_\%A\_\%a.out}
\CommentTok{\#SBATCH {-}{-}error=logs/fastp\_array\_\%A\_\%a.err}
\CommentTok{\#SBATCH {-}{-}array=1{-}4}
\CommentTok{\#SBATCH {-}{-}ntasks=1}
\CommentTok{\#SBATCH {-}{-}cpus{-}per{-}task=2}
\CommentTok{\#SBATCH {-}{-}mem=5G}
\CommentTok{\#SBATCH {-}{-}time=01:00:00}

\CommentTok{\# Make conda available in the job, and activate the environment with fastp}
\BuiltInTok{source}\NormalTok{ \textasciitilde{}/.bashrc}
\ExtensionTok{conda}\NormalTok{ activate fastp\_1.4.0}

\CommentTok{\# Take the line with the number of this task from samples.txt}
\VariableTok{sample}\OperatorTok{=}\VariableTok{$(}\FunctionTok{sed} \AttributeTok{{-}n} \StringTok{"}\VariableTok{$\{SLURM\_ARRAY\_TASK\_ID\}}\StringTok{p"}\NormalTok{ samples.txt}\VariableTok{)}
\BuiltInTok{echo} \StringTok{"Start fastp for: }\VariableTok{$sample}\StringTok{"}

\ExtensionTok{fastp} \DataTypeTok{\textbackslash{}}
    \AttributeTok{{-}i}\NormalTok{ data/seq\_project/}\PreprocessorTok{*}\NormalTok{/}\VariableTok{$\{sample\}}\NormalTok{\_R1.fastq.gz }\DataTypeTok{\textbackslash{}}
    \AttributeTok{{-}I}\NormalTok{ data/seq\_project/}\PreprocessorTok{*}\NormalTok{/}\VariableTok{$\{sample\}}\NormalTok{\_R2.fastq.gz }\DataTypeTok{\textbackslash{}}
    \AttributeTok{{-}o}\NormalTok{ results/fastp\_array/}\VariableTok{$\{sample\}}\NormalTok{\_R1.trimmed.fastq.gz }\DataTypeTok{\textbackslash{}}
    \AttributeTok{{-}O}\NormalTok{ results/fastp\_array/}\VariableTok{$\{sample\}}\NormalTok{\_R2.trimmed.fastq.gz }\DataTypeTok{\textbackslash{}}
    \AttributeTok{{-}h}\NormalTok{ results/fastp\_array/}\VariableTok{$\{sample\}}\NormalTok{.html }\DataTypeTok{\textbackslash{}}
    \AttributeTok{{-}j}\NormalTok{ results/fastp\_array/}\VariableTok{$\{sample\}}\NormalTok{.json }\DataTypeTok{\textbackslash{}}
    \AttributeTok{{-}{-}average\_qual}\NormalTok{ 20 }\DataTypeTok{\textbackslash{}}
    \AttributeTok{{-}{-}thread} \VariableTok{$\{SLURM\_CPUS\_PER\_TASK\}}
\end{Highlighting}
\end{Shaded}

\begin{itemize}
\tightlist
\item
  \texttt{-\/-array=1-4}: one task per line in \texttt{samples.txt}.
  Check the number with \texttt{wc\ -l\ samples.txt}
\item
  \texttt{sed\ -n\ "\$\{SLURM\_ARRAY\_TASK\_ID\}p"\ samples.txt}: prints
  only line number \texttt{SLURM\_ARRAY\_TASK\_ID} of the file. Task 1
  gets the first sample, task 2 the second, and so on
\item
  \texttt{-\/-thread\ \$\{SLURM\_CPUS\_PER\_TASK\}}: SLURM stores the
  number of CPUs you requested with \texttt{-\/-cpus-per-task} in this
  variable. fastp now always uses as many CPUs as the job has, also if
  you change the number later
\item
  Removed: \texttt{DATA\_DIR}, \texttt{OUT\_DIR}, the check on the input
  files (check the inputs before you submit instead), the comments at
  the end of the lines and \texttt{echo\ "Done"}
\item
  Left out on purpose: the quality trimming. The prompt asked for it,
  but not for a specific method. fastp has three ways to trim reads by
  quality (\texttt{-\/-cut\_front}, \texttt{-\/-cut\_tail} and
  \texttt{-\/-cut\_right}, see the answer to the worked example and
  \texttt{fastp\ -\/-help}). Read about them first, decide which one
  fits your data, and only then add it, with a note in your notebook on
  why you chose it
\item
  Added: \texttt{conda\ activate}, \texttt{samples.txt}, and our own
  names for the reports
\item
  Make the output folder before you submit:
  \texttt{mkdir\ -p\ logs\ results/fastp\_array}
\end{itemize}

Before you run it on all samples: set \texttt{-\/-array=1} to run only
the first sample, and check the output and the fastp report.
\end{quote}

\end{tcolorbox}

\section{When your code does not
work}\label{when-your-code-does-not-work}

When your code crashes or gives strange results, it is tempting to paste
everything into an LLM and ask it to fix it. I suggest a different
order:

\begin{enumerate}
\def\labelenumi{\arabic{enumi}.}
\tightlist
\item
  \textbf{Read the error message.} As on day 1: which command complains,
  and about which file? Often, the message already tells you what is
  wrong
\item
  \textbf{Narrow it down.} Bugs are almost always small and in one
  place. Run the steps one by one. Put \texttt{echo} in front of a
  command to see what bash actually runs, or \texttt{echo\ \$sample} to
  see what is in a variable
\item
  \textbf{Then ask a specific question.} Give the LLM the one command
  that fails, the exact error message and a few lines of the input
  (replace anything sensitive, see above). Ask what the error means, not
  for a new script
\end{enumerate}

A rewritten script may run without errors, but you then have to check
all of it again, and you have learned nothing about why the first
version failed.

\section{Further reading}\label{further-reading}

\begin{itemize}
\tightlist
\item
  \href{https://scienceparkstudygroup.github.io/ibed-bioinformatics-page/source/core_tools/bash-for-loops.html}{On
  GenAI usage}, at the end of the IBED for-loops page: what to look out
  for in loops written by an LLM
\item
  The UvA's \href{https://edu.nl/gakhm}{AI literacy course for
  students}: what generative AI is, how it works and what its limits are
  (about 45 to 60 minutes)
\item
  \href{https://www.uva.nl/en/about-the-uva/about-the-university/ai/ai-in-education/uva-ai-chat/uva-ai-chat.html}{UvA
  AI Chat} and the UvA
  \href{https://student.uva.nl/en/information/ai-tools-and-your-studies}{rules
  for students}
\end{itemize}




\end{document}
